% from main.tex
%%%%%%%%%%  scan idx 162  |  folio 149  %%%%%%%%%%
Supposons que \(E\) soit en outre de tor-dimension finie (condition réalisée si \(S\) est de type fini). Alors, si \(F\) est à cohomologie bornée (resp. de tor-dimension finie), il en est de même de \(\RHom(E,F)\). Plus précisément, soient \(a,b,m,n\) des entiers tels que \(\tamp(E)\subset[a,b]\) (ou, ce qui revient au même, d'après (5.8), \(\pamp(E)\subset[a,b]\)) et \(\Ho^i(F)=0\) pour \(i\notin[m,n]\) (resp. \(\tamp(F)\subset[m,n]\)). Alors on a \(\Ho^i(\RHom(E,F))=0\) pour \(i\notin[m-b,n-a]\) (resp. \(\tamp(\RHom(E,F))\subset[m-b,n-a]\)).

\emph{Preuve.} L'assertion a) est locale, donc on peut supposer \(E\) strictement parfait, et l'on gagne par dévissage sur \(E\) (en se ramenant à \(E=A\)). Pour l'assertion b), on peut supposer \(E\) strictement parfait tel que \(E^i=0\) pour \(i\notin[a,b]\), et \(F^i=0\) pour \(i\notin[m,n]\) (resp. \(F^i\) plat pour tout \(i\) et \(F^i=0\) pour \(i\notin[m,n]\)). Il suffit alors d'appliquer le lemme suivant :

\textbf{Lemme 7.1.1.} Soit \(E\in\Ob\Dpx^-(S)\), \(F\in\Ob\Dpx^+(S)\). Si \(E\) est strictement pseudo-cohérent (i.e. (2.1) à composantes facteurs directs de Modules libres de type fini), la flèche canonique de \(\Dpx^+(S)\)
\[
  \Hom^\bullet(E,F)\longrightarrow\RHom(E,F)
\]
est un isomorphisme.

\emph{Preuve.} La suite spectrale
\[
  E_2^{pq}=\Ho^p(\Ext^q(E,F)^\bullet)\Longrightarrow \Ext^*(E,F),
\]
où
\[
  \Ext^q(E,F)^i=\coprod_{i_2-i_1=i}\Ext^q(E^{i_1},F^{i_2}),
\]
dégénère puisque \(\Ext^q(E,F)^\bullet=0\) pour tout \(q>0\) (1.3.1). On en déduit des
